\chapter{Planteamiento del problema y objetivos}
\label{cap:planteamiento}

\section{Planteamiento del problema}
El despliegue de tecnolog\'ias de captaci\'on pasiva de agua atmosf\'erica en regiones \'aridas y semi\'aridas costeras se enfrenta a un doble desaf\'io f\'isico: la necesidad de maximizar el volumen de agua condensada y recolectada durante los episodios de niebla y alta humedad, frente a la degradaci\'on de la membrana colectora provocada por la deposici\'on y cristalizaci\'on de aerosoles marinos. Las estrategias de mitigaci\'on basadas en temporizadores operan sin retroalimentaci\'on del estado f\'isico de la superficie ni de las condiciones del entorno. Esto puede producir ciclos t\'ermicos innecesarios que consumen energ\'ia de reserva, acumulan fatiga t\'ermica y no aseguran la regeneraci\'on de la capacidad de flujo cuando la incrustaci\'on salina crece m\'as r\'apido de lo que el programa fijo puede retirar.

Surge as\'i la necesidad de formalizar un esquema de control reactivo, computable en hardware embebido de bajo consumo, que tome decisiones a partir de par\'ametros f\'isicos medibles (temperatura de la matriz, humedad relativa ambiental, impedancia el\'ectrica superficial y estado de carga de la bater\'ia), jerarquizando la integridad de los componentes y la preservaci\'on del equilibrio ecol\'ogico.

\section{Pregunta de investigaci\'on}
\label{sec:pregunta}
Con base en lo anterior, el proyecto define la siguiente pregunta central de investigaci\'on:

\begin{quote}
\itshape
?`En qu\'e medida un sistema de control reactivo basado en reglas, como primera etapa de una arquitectura neuro-simb\'olica, optimiza la eficiencia de captaci\'on pasiva de agua atmosf\'erica y previene la degradaci\'on por incrustaci\'on salina en un colector de niebla ubicado en una zona semi\'arida y costera como Cabo de la Vela (La Guajira), en comparaci\'on con un sistema convencional de temporizaci\'on peri\'odica?
\end{quote}

Para que la pregunta sea contrastable, sus t\'erminos se definen operacionalmente sin cambiar su enunciado. La \emph{eficiencia de captaci\'on} se mide con dos variables: la recuperaci\'on post-pulso y el volumen de agua colectado por la red en la corrida, en litros o en $\text{L/m}^2\text{/d\'ia}$. La \emph{recuperaci\'on post-pulso} de un pulso es la conductancia relativa que queda al terminarlo, $R = \min\left(100,\; 100\cdot Z_{\text{nom}}/Z_{\text{tras pulso}}\right)\,\%$, con $Z_{\text{nom}} = 1000\;\Omega$, y se promedia sobre todos los pulsos de cada modo de control. La \emph{degradaci\'on} es la p\'erdida acumulada de integridad de la membrana, $100 - I$, con $I$ en porcentaje. La \emph{carga salina} es el aumento de impedancia que deposita cada litro captado por metro cuadrado y se expresa en $\Omega$ por $\text{L/m}^2$, porque el modelo usa la impedancia como sustituto el\'ectrico de la sal acumulada, sin medir su masa. Por \'ultimo, la \emph{incrustaci\'on diaria}, en $\Omega\text{/d\'ia}$, es el producto de la carga salina por el rendimiento diario de agua: el aumento de impedancia que recibir\'ia la membrana en un d\'ia de operaci\'on.

De esta pregunta general se derivan cuatro subpreguntas operativas:
\begin{enumerate}[label=\alph*)]
  \item ?`C\'omo incide el control basado en reglas en la reducci\'on del desgaste y la degradaci\'on acumulada de la membrana frente al temporizador peri\'odico?
  \item ?`A partir de qu\'e nivel de carga e incrustaci\'on salina diaria el control basado en reglas logra una ventaja efectiva en la recuperaci\'on post-pulso frente al temporizador?
  \item ?`En qu\'e medida garantiza el control basado en reglas la seguridad t\'ermica, la viabilidad energ\'etica y la no violaci\'on de restricciones ecol\'ogicas en comparaci\'on con el control tradicional?
  \item ?`Cu\'al es el impacto y dimensionamiento requerido en t\'erminos de personas abastecidas y equivalencia con la entrega de agua asistencial (ICBF) en distintos entornos de niebla y salinidad?
\end{enumerate}

\section{Objetivos}

\subsection{Objetivo general}
Evaluar mediante simulaci\'on computacional comparativa el desempe\~no de un motor de control reactivo basado en reglas, como primera etapa de una arquitectura neuro-simb\'olica, frente a un sistema de temporizaci\'on peri\'odica, en t\'erminos de eficiencia de captaci\'on pasiva de agua atmosf\'erica, mitigaci\'on de la degradaci\'on por incrustaci\'on salina, seguridad operativa e impacto potencial en zonas costeras semi\'aridas como Cabo de la Vela.

\subsection{Objetivos espec\'ificos}
\begin{enumerate}[leftmargin=*]
  \item Dise\~nar e implementar un motor de seis reglas deterministas jerarquizadas seg\'un criterios de seguridad f\'isica, conservaci\'on ecol\'ogica, protecci\'on energ\'etica y regeneraci\'on por impedancia superficial.
  \item Desarrollar un modelo f\'isico-ambiental y un simulador continuo que reproduzca la din\'amica clim\'atica nocturna/diurna de Cabo de la Vela, la formaci\'on de pel\'icula de agua, la acumulaci\'on salina y la degradaci\'on estructural de la membrana.
  \item Realizar barridos de salinidad, an\'alisis de sensibilidad global, propagaci\'on de incertidumbre y pruebas de robustez del cruce de recuperaci\'on y de la reducci\'on del desgaste ante variaciones en los par\'ametros mec\'anicos y ambientales.
  \item Calibrar el flujo de niebla del modelo contra los rendimientos documentados en la literatura internacional y estimar el dimensionamiento f\'isico del sistema y su equivalencia frente a las entregas de agua por carrotanque reportadas en el territorio.
\end{enumerate}

\section{Justificaci\'on multidimensional}

\subsection{Perspectiva t\'ecnica}
Desde la \'optica del dise\~no de sistemas de control para dispositivos desatendidos, un conjunto peque\~no de reglas deterministas tiene tiempos de ejecuci\'on acotados y un comportamiento auditable: cada decisi\'on puede rastrearse hasta la regla que la produjo. El modelo supone un consumo base de $0.6\;\text{W}$ para la electr\'onica y los sensores del control por reglas (\cref{anexo:trazabilidad}), un valor de orden de magnitud que deber\'a medirse en hardware. Esa base simb\'olica es, adem\'as, el componente que permitir\'ia incorporar despu\'es modelos neuronales predictivos sin renunciar a los l\'imites de seguridad.

\subsection{Perspectiva social y humana}
El acceso al agua potable es un derecho humano reconocido internacionalmente \parencite{onu2010agua}, y en La Guajira su garant\'ia fue objeto de la Sentencia T-302 de 2017 \parencite{corte2017t302}. Dimensionar sistemas aut\'onomos de captaci\'on pasiva con est\'andares expl\'icitos de la OMS ($20\;\text{L/persona/d\'ia}$) \parencite{howard2020domestic} y de Sphere ($15\;\text{L/persona/d\'ia}$) \parencite{sphere2018handbook} permite traducir hip\'otesis tecnol\'ogicas en par\'ametros tangibles de bienestar humano.

\subsection{Perspectiva \'etica y cultural}
Cabo de la Vela es territorio ancestral de la naci\'on Wayuu. Cualquier innovaci\'on tecnol\'ogica debe respetar el r\'egimen jur\'idico de consulta previa, libre e informada, que en Colombia se deriva del Convenio 169 de la Organizaci\'on Internacional del Trabajo, aprobado mediante la Ley 21 de 1991 \parencite{ley21de1991}, as\'i como las pautas de gobernanza territorial comunitaria. En consonancia con la Sentencia T-302 de 2017 \parencite{corte2017t302}, que reconoce el acceso al agua potable como un derecho fundamental interdependiente con la salud y la alimentaci\'on de la ni\~nez wayuu, el equipo asume como postura de dise\~no que el agua no debe tratarse \'unicamente como un recurso termodin\'amico o insumo econ\'omico, sino como un bien com\'un ligado al equilibrio del territorio. Por ello, el artefacto incorpora expl\'icitamente el l\'imite de extracci\'on ecol\'ogica y la pausa en sequ\'ia para no competir con el microclima circundante. Esta justificaci\'on abre una pregunta emp\'irica para el Ciclo 3: si la potencia el\'ectrica media del nodo en hardware embebido se mantiene en torno a los $0.6\;\text{W}$ supuestos y cu\'antos litros diarios por persona deja de entregar la red cuando la comunidad ajusta el tope de extracci\'on y el umbral de humedad de la pausa ecol\'ogica.

Delimitados el problema, la pregunta y sus definiciones, el cap\'itulo siguiente sit\'ua la propuesta en la literatura: qu\'e se sabe de la captaci\'on de niebla, de la incrustaci\'on en membranas y del control de limpieza basado en condici\'on, y d\'onde se ubica el aporte de este trabajo.
